% ===================================================================
%  丛书共用前置件
%
%  目录、图表目录、丛书页三者，各卷写法完全相同，只差一处高亮。
%  在这里各做成一个命令，各卷 frontmatter 只调用，不再重抄版面代码。
%
%  依赖：series-base.tex 已载入。
% ===================================================================

% ---- 目录与图表目录 -------------------------------------------------
% 用法：\tocandfigs
% 说明：\tableofcontents 自己会开新页，故用底层命令手写标题。
\newcommand{\tocandfigs}{%
  \clearpage
  \chapter*{目\quad 录}%
  \phantomsection
  \addcontentsline{toc}{chapter}{目\quad 录}%
  \markboth{目\quad 录}{目\quad 录}%
  \vspace{2mm}%
  {\linespread{1.16}\selectfont
   \csname @starttoc\endcsname{toc}%
  }%
  \clearpage
  \chapter*{图表目录}%
  \phantomsection
  \addcontentsline{toc}{chapter}{图表目录}%
  \markboth{图表目录}{图表目录}%
  \vspace{2mm}%
  {\linespread{1.16}\selectfont
   {\noindent\sffamily\normalsize\bfseries 图\quad 版\par}%
   \vspace{2mm}%
   \csname @starttoc\endcsname{lof}%
   \vspace{8mm}%
   {\noindent\sffamily\normalsize\bfseries 表\quad 格\par}%
   \vspace{2mm}%
   \csname @starttoc\endcsname{lot}%
  }%
}

% ---- 丛书页 ---------------------------------------------------------
% 用法：\seriespage{vii}
% 参数为本书的卷次（小写罗马数字），该行加粗；其余各行照排。
\newcommand{\seriespage}[1]{%
  \clearpage
  \thispagestyle{empty}%
  \begin{center}
  \vspace*{12mm}
  {\sffamily\small 远东异端史料与文明起源重构丛书}\\[3mm]
  {\sffamily\footnotesize\color{pencil}%
   Восточные еретические источники и реконструкция происхождения цивилизаций}\\[10mm]
  {\sffamily\small 主编\quad А.\,В.\,巴兰尼科夫}
  \end{center}
  \vspace{6mm}
  {\footnotesize\linespread{1.5}\selectfont
   \begin{list}{}{\leftmargin=8mm \itemindent=-8mm \topsep=0pt \parsep=0pt}
   \seriesrow{i}{《月之都——〈广寒宫残简〉辑佚与考释》}{巴兰尼科夫 辑校}{#1}
   \seriesrow{ii}{西周“天车”铭文与星际物流}{库兹涅佐夫}{#1}
   \seriesrow{iii}{秦俑坑第 4 号坑未被发掘的部分}{谢列兹尼奥娃}{#1}
   \seriesrow{iv}{汉代地动仪作为引力波天线之功能}{巴兰尼科夫}{#1}
   \seriesrow{v}{三国时期的电磁抛射技术与后勤}{彼得罗夫}{#1}
   \seriesrow{vi}{敦煌卷子中的密钥学残卷}{伊万年科}{#1}
   \seriesrow{vii}{战争的雷霆——记一种可能的中国历史}{巴兰尼科夫}{#1}
   \seriesrow{viii}{唐长安城反物质管网复原}{巴兰尼科夫}{#1}
   \seriesrow{ix}{明代锦衣卫外骨骼兵人体质档案}{张建国}{#1}
   \seriesrow{x}{郑和舰队南极补给点考}{巴兰尼科夫}{#1}
   \seriesrow{xi}{太医院基因工坊编年}{巴兰尼科夫}{#1}
   \seriesrow{xii}{记忆消除场的文献学后果}{索科洛夫}{#1}
   \seriesrow{xiii}{月之都——记一个不想变化的文明}{[日]\,太田顺也}{#1}
   \end{list}}
  \vspace{10mm}
  {\footnotesize\color{pencil}
   本丛书所收，均为对远东地区「异端史料」之重新整理与研究。\\
   各卷可独立阅读，互不倚赖；引用遵循各卷自行之体例。\\[2mm]
   全书共十三卷；第十三卷卷次暂定。}
  \clearpage
}

% 单行。第 4 个参数为本书卷次，相符则整行加粗。
\makeatletter
\newcommand{\seriesrow}[4]{%
  \def\@tmp{#4}%
  \def\@mine{#1}%
  \ifx\@tmp\@mine
    \item \textbf{\textsc{#1}}\hspace{3mm}\textbf{#2}\hfill\textbf{#3}%
  \else
    \item \textsc{#1}\hspace{4mm}#2\hfill #3%
  \fi
}
\makeatother

% ---- 版权页 ---------------------------------------------------------
% 用法：\copyrightpage{<版权页表格行>}
% 参数为 tabular 的行体，用 \cfrow{项目}{内容} 逐行给出。
% 版式（minipage、页脚、法律声明）一律由本命令提供，各卷只填内容。
\newcommand{\cfrow}[2]{#1 & #2\\}
\newcommand{\copyrightpage}[1]{%
  \clearpage
  \thispagestyle{empty}%
  \noindent\begin{minipage}[t][\dimexpr\textheight-24pt\relax][s]{\textwidth}
  \footnotesize
  \vspace*{4mm}
  \raggedright
  \renewcommand{\arraystretch}{1.05}%
  \begin{tabular}{@{}p{17mm}p{0.80\textwidth}@{}}
  #1
  \end{tabular}
  \vfill
  {\scriptsize\color{pencil}
  版权所有\ \ 侵权必究。未经出版者书面许可，不得以任何方式复制或抄袭本书之部分或全部内容。\\
  本书如有印装质量问题，请与本社读者服务部联系调换。}
  \end{minipage}
  \clearpage
}
